\section{第39届复赛第7题：高能铅核碰撞夸克物质状态方程}
\subsection{题干}
高能铅核碰撞形成的高温球形“火球”近似由静止质量很小、速度接近光速的夸克组成。夸克为无相互作用的经典粒子，速度分布各向同性，能量分布比例为 $e^{-E/(k_BT)}4\pi p^2\mathrm dp$。求：（1）$P$ 与能量密度 $u$ 的状态方程；（2）$P,n,T$ 的状态方程；（3）$C_V$ 与热容比 $\gamma$；（4）半径 $3.0\times10^{-15}\,\mathrm m$、温度从 $400\,\mathrm{MeV}/k_B$ 降至 $150\,\mathrm{MeV}/k_B$ 时的火球半径。
\subsection{参考解答整理}
\paragraph{（1）器壁动量通量。} 对速度组 $i$，近光速近似给出 $\boldsymbol p_i\cdot\boldsymbol v_i\simeq E_i$。由各向同性，器壁压强
\[P=\frac13\sum_i n_i\boldsymbol p_i\cdot\boldsymbol v_i=\boxed{\frac13u}.\]
\begin{figure}[htbp]\centering\begin{tikzpicture}[>=Stealth,scale=0.78]
  \draw[very thick] (3.3,-1.8) -- (3.3,2.4);
  \draw[->,thick] (-2.8,0) -- (4.4,0) node[right] {$x$};
  \draw[<->] (3.3,-1.08) -- (-1.25,-1.08) node[midway,below] {$v_{ix}\Delta t$};
  \draw[->,thick] (-2.7,2.05) -- (3.25,0.62);
  \draw[->,thick] (-2.45,1.55) -- (3.25,0.36);
  \draw[->,thick] (-2.15,1.1) -- (3.25,0.1);
  \draw[->,thick] (-2.62,-1.52) -- (3.25,-0.25);
  \draw[->,thick] (-2.32,-1.1) -- (3.25,-0.05);
  \draw[->,thick] (-2.0,-0.7) -- (3.25,0.18);
  \draw[densely dashed] (-1.25,-1.42) -- (-1.25,0.32) -- (0.25,0.83) -- (0.25,-1.42) -- cycle;
  \node at (-0.5,-0.28) {$\Delta S$};
  \node[below] at (3.3,-1.8) {解题图7a：速度组粒子在 $\Delta t$ 内撞击面元};
\end{tikzpicture}
\caption{依据答案解题图7a重绘的器壁面元与粒子束。}\end{figure}
\begin{figure}[htbp]\centering\begin{tikzpicture}[>=Stealth,scale=1.05]
  \draw[very thick] (2.4,-1.8) -- (2.4,1.8);
  \draw[dashed] (-0.65,0) -- (2.4,0);
  \draw[->,thick] (-0.3,-1.4) -- (2.33,-0.25);
  \draw[->,thick] (-0.7,-0.82) -- (2.33,0.48);
  \draw[->,thick] (0.05,1.1) -- (2.33,0.1);
  \draw[->] (1.72,0) arc[start angle=180,end angle=157,radius=.68] node[midway,left] {$\theta$};
  \draw[<->] (2.62,-0.25) -- (2.62,0.48) node[midway,right] {$\mathrm dS$};
  \node[below] at (0.9,-1.95) {解题图7b：粒子入射角与器壁面元};
\end{tikzpicture}
\caption{依据答案解题图7b重绘的入射角积分几何。}\end{figure}
\paragraph{（2）平均能量与状态方程。} 用 $E\simeq pc$ 代入题设分布，积分比值给出 $\bar E=3k_BT$。故
\[\boxed{P=\frac13n\bar E=nk_BT}.\]
\paragraph{（3）热容。} 一摩尔内能 $U_m=N_A\bar E=3RT$，所以
\[\boxed{C_V=3R,\qquad C_P=4R,\qquad \gamma=\frac43}.\]
\paragraph{（4）绝热膨胀。} 准静态绝热过程 $TV^{\gamma-1}=\mathrm{const}$，代入 $\gamma=4/3$ 和 $V=4\pi r^3/3$ 得 $Tr=\mathrm{const}$，故
\[\boxed{r=\frac{400}{150}\times3.0\times10^{-15}\,\mathrm m=8.0\times10^{-15}\,\mathrm m}.\]
\subsection{转写依据}
题干依据 \texttt{CPHO39-Q07.md} 第4页，答案依据第27至29页与 \texttt{CPHO39-Q07\_formula\_check\_v1.csv}。两幅图分别依据答案解题图7a、7b重绘。本题四问分值依次为11、10、10、9分。
